% EM structure and teaching revision 20260928
\subsection{无限网络的递推、选根与收敛}
自相似结构允许用整体的等效量表示去掉首段后剩下的部分。设每段先串联一个 $R$，再把一个 $R$ 对地分流，余下网络等效仍为 $X$，则
\[
 X=R+\frac{RX}{R+X},\qquad X^2-RX-R^2=0.
\]
正电阻网络取 $X=(1+\sqrt5)R/2$。负根虽满足代数方程，却不满足被动端口的非负耗散条件。
更可靠的定义来自有限截断：从某个非负末端电阻 $X_0$ 出发，令 $X_{n+1}=R+RX_n/(R+X_n)$。一步后有 $R\le X_n\le2R$，且该区间内导数
$f'(X)=R^2/(R+X)^2\le1/4$，所以相邻近似间的差不断缩小，极限确为所选正根。
有多个端口时，一个标量等效电阻往往不够。应保留端口电势和注入电流的线性关系，或者利用对称性把端口激励分解成独立模式。未经证明就把三端网络替换成两端电阻，会丢失约束，使看似正确的递推求出错误结果。

